% ===========================================================================
%  Elementos pre-textuais externos (capa, folha de rosto, aprovacao, etc.)
% ===========================================================================
\thispagestyle{empty}

% ------------------------------- CAPA --------------------------------------
\begin{center}
  {\bfseries\large INSTITUTO FEDERAL FLUMINENSE}\\[2pt]
  {\bfseries CURSO DE ENGENHARIA [VALIDAR COM O AUTOR]}\\[1pt]
  \rule{\textwidth}{0.4pt}
\end{center}

\vspace{2.6cm}

\begin{center}
  {\bfseries\Large WILLIAM DA SILVA VIANNA}
\end{center}

\vspace{3.2cm}

\begin{center}
  {\bfseries\Large MODELAGEM 3D PROCEDIMENTAL DO LOTUS ELISE 111S EM BLENDER\\[6pt]
   AUTOMATIZADA POR MODEL CONTEXT PROTOCOL (MCP)}
\end{center}

\vfill

\begin{center}
  {\bfseries CAMPOS DOS GOYTACAZES -- RJ}\\[2pt]
  {\bfseries 2026}
\end{center}

\clearpage

% --------------------------- FOLHA DE ROSTO ---------------------------------
\thispagestyle{empty}
\begin{center}
  {\bfseries\large WILLIAM DA SILVA VIANNA}
\end{center}

\vspace{3.0cm}

\begin{center}
  {\bfseries\Large MODELAGEM 3D PROCEDIMENTAL DO LOTUS ELISE 111S EM BLENDER\\[6pt]
   AUTOMATIZADA POR MODEL CONTEXT PROTOCOL (MCP)}
\end{center}

\vspace{1.5cm}

\hfill
\begin{minipage}{0.55\textwidth}
  \SingleSpacing
  \footnotesize
  \raggedright
  Monografia apresentada ao Curso de Engenharia [VALIDAR COM O AUTOR] do Instituto Federal
  Fluminense, como requisito parcial para obten\c{c}\~ao do t\'itulo de Engenheiro
  [VALIDAR COM O AUTOR].\\[8pt]
  \'Area de concentra\c{c}\~ao: Computa\c{c}\~ao Gr\'afica e Automa\c{c}\~ao de Pipelines de
  Modelagem.\\[8pt]
  Autor e orientador: William da Silva Vianna (o pr\'oprio autor), Instituto Federal
  Fluminense.
\end{minipage}

\vfill

\begin{center}
  {\bfseries CAMPOS DOS GOYTACAZES -- RJ}\\[2pt]
  {\bfseries 2026}
\end{center}

\clearpage

% ------------------------- FOLHA DE APROVACAO -------------------------------
\thispagestyle{empty}
\begin{center}
  {\bfseries\large WILLIAM DA SILVA VIANNA}
\end{center}
\vspace{2.0cm}
\begin{center}
  {\bfseries\Large MODELAGEM 3D PROCEDIMENTAL DO LOTUS ELISE 111S EM BLENDER\\
   AUTOMATIZADA POR MODEL CONTEXT PROTOCOL (MCP)}
\end{center}
\vspace{1.6cm}

\hfill
\begin{minipage}{0.62\textwidth}
  \SingleSpacing
  \footnotesize
  \raggedright
  Monografia apresentada ao Curso de Engenharia [VALIDAR COM O AUTOR] do Instituto Federal
  Fluminense, como requisito parcial para obten\c{c}\~ao do t\'itulo de Engenheiro
  [VALIDAR COM O AUTOR].
\end{minipage}

\vspace{1.4cm}

\begin{flushright}
  \footnotesize
  Aprovada em: \rule{3cm}{0.4pt}
\end{flushright}

\vspace{1.6cm}

\begin{center}
  \footnotesize
  \begin{tabular}{>{\centering\arraybackslash}p{7.6cm}}
    \hline
    William da Silva Vianna -- Autor e Orientador\\[-2pt]
    Instituto Federal Fluminense
  \end{tabular}

  \vspace{1.4cm}

  \begin{tabular}{>{\centering\arraybackslash}p{7.6cm}}
    \hline
    Prof.\ [VALIDAR COM O AUTOR] -- Membro da banca\\[-2pt]
    [VALIDAR COM O AUTOR]
  \end{tabular}

  \vspace{1.4cm}

  \begin{tabular}{>{\centering\arraybackslash}p{7.6cm}}
    \hline
    Prof.\ [VALIDAR COM O AUTOR] -- Membro da banca\\[-2pt]
    [VALIDAR COM O AUTOR]
  \end{tabular}
\end{center}

\clearpage

% ------------------------------ DEDICATORIA ---------------------------------
\thispagestyle{empty}
\vspace*{\fill}
\begin{flushright}
  \begin{minipage}{0.5\textwidth}
    \SingleSpacing
    \itshape
    Aos que entendem a modelagem como engenharia: onde cada v\'ertice \'e uma decis\~ao
    documentada e cada superf\'icie, uma hip\'otese test\'avel.
  \end{minipage}
\end{flushright}
\vspace*{\fill}
\clearpage

% ---------------------------- AGRADECIMENTOS --------------------------------
\thispagestyle{empty}
\begin{center}
  {\bfseries AGRADECIMENTOS}
\end{center}
\vspace{1.0cm}

\noindent
Agrade\c{c}o ao Instituto Federal Fluminense pela forma\c{c}\~ao e pela infraestrutura que
tornaram este trabalho poss\'ivel, e ao corpo docente do Curso de Engenharia pelo rigor que
moldou a maneira como este projeto foi conduzido.

\noindent
Registro que este trabalho foi conduzido com a orienta\c{c}\~ao do pr\'oprio autor, William da
Silva Vianna, que acumulou as fun\c{c}\~oes de orientando e orientador ao longo de todo o
desenvolvimento --- arranjo que, se por um lado concentrou as decis\~oes t\'ecnicas, por outro
exigiu que cada uma delas fosse registrada de forma expl\'icita neste documento.

\noindent
A Victor Vianna, meu sobrinho, pela colabora\c{c}\~ao com a ideia da modelagem. A sugest\~ao de
tomar o Lotus Elise 111S como objeto de estudo foi dele, e sem ela este trabalho n\~ao
existiria na forma que assumiu.

\noindent
\'A comunidade de software livre, em especial aos mantenedores do Blender, cuja
documenta\c{c}\~ao p\'ublica e c\'odigo-fonte aberto permitiram que cada decis\~ao deste
trabalho fosse verificada na fonte.

\vspace{1.0cm}
\clearpage

% -------------------------------- EPIGRAFE ----------------------------------
\thispagestyle{empty}
\vspace*{\fill}
\begin{flushright}
  \begin{minipage}{0.55\textwidth}
    \SingleSpacing
    \itshape
    ``O que n\~ao est\'a no script n\~ao existe.''
  \end{minipage}
\end{flushright}
\vspace*{\fill}
\clearpage
